\chapter{今後の展望}

\section{アルゴリズムの拡張}

\subsection{動的環境への対応}

本研究では、事前に生成された静的な地形マップを用いて経路計画を行った。しかし、実環境では、天候の変化（雨による地面の軟化、雪の堆積）や、新たな障害物の出現により、地形情報が動的に変化する。今後の展望として、ロボットが走行中にセンサーから得られる情報を用いて、地形マップをオンラインで更新する手法の開発が挙げられる。

動的地形更新には、RTAB-MapのSLAM機能を活用できる。RTAB-Mapは、リアルタイムで環境地図を更新する機能を持つため、ロボットが走行しながら新たに観測した地形情報を地形マップに反映できる。更新された地形マップに基づいて、Field D* Hybridの再計画機能を用いて、経路を動的に更新する。これにより、環境変化に適応した経路計画が可能となり、より頑健なシステムを実現できる。

\subsection{マルチモーダル地形評価}

本研究では、標高差、傾斜角、表面粗さの3要素から地形複雑度を計算した。今後は、RGB-Dカメラから得られる視覚情報（テクスチャ、色）を統合することで、より詳細な地形評価が可能となる。例えば、植生の密度、水たまりの有無、岩の種類などを識別し、それぞれに対する走行適性を評価できる。このマルチモーダルな地形評価により、より精密な経路計画が実現できる。

視覚情報の統合には、深層学習を用いた地形分類が有効である。事前に学習したモデルを用いて、カメラ画像から地形タイプ（草地、砂地、岩石など）を分類し、各タイプに対する地形コストを設定する。これにより、点群データだけでは識別困難な地形特性も考慮できる。

\subsection{機械学習による地形予測}

さらに、機械学習を用いた地形予測も有望である。ロボットが走行中に収集したデータ（地形情報、消費エネルギー、走行速度）を蓄積し、地形とロボット性能の関係を学習する。これにより、各ロボットの個別特性（車輪の種類、重量、モーター性能）に適応した地形評価が可能となる。

例えば、クローラー型ロボットと車輪型ロボットでは、同じ地形でも走行適性が異なる。クローラー型は砂地や泥濘地での走行に優れるが、硬い路面では効率が劣る。車輪型はその逆の特性を持つ。機械学習により、各ロボットに最適化された地形コストを学習することで、より効率的な経路計画が実現できる。

\section{実環境実験}

\subsection{フィールド実証実験の必要性}

本研究では、シミュレーション環境において提案手法の有効性を検証した。今後は、実環境におけるフィールド実証実験が必要である。実環境では、センサーノイズ、通信遅延、予期せぬ障害物など、シミュレーションでは再現が困難な要因が存在する。実機ロボットを用いてフィールド実験を行い、提案手法が実環境でも有効に機能することを確認する必要がある。

フィールド実験では、海岸、山地、農地などの実際の不整地環境において、ロボットを走行させる。RTAB-Mapによる環境地図生成、地形情報抽出、経路計画、経路追従制御の一連のシステムを統合し、完全自律での移動を実現する。実験を通じて、システムの頑健性、計算効率、安全性を評価し、実用化に向けた知見を得る。

\subsection{多様な環境での評価}

また、海岸、山地、森林など、様々な地形条件において提案手法の性能を評価することも重要である。本研究では、丘陵地形を主な対象としたが、実用化に向けては、多様な環境での動作確認が必要である。

さらに、季節や天候の変化（晴天、雨天、雪）による地形特性の変化にも対応できることを確認する必要がある。雨天時は地面が滑りやすくなり、地形コストが増加する。雪が積もると、地形の凹凸が埋まり、地形情報が変化する。これらの変化に対して、ロボットがどのように経路を調整するかを評価する。

長距離走行試験により、バッテリー消費や経路計画の安定性を評価することも今後の課題である。数km以上の長距離移動において、累積誤差やメモリ使用量の増加がシステムに与える影響を検証する。

\section{他分野への応用}

\subsection{農業ロボットへの応用}

提案手法は、農業ロボットへの応用が期待できる。農業ロボットは、畑の畝や溝を避けながら移動し、作物を傷つけずに作業する必要がある。TA*を用いることで、畝を地形の悪い領域として扱い、畝を避けた経路を生成できる。また、土壌の硬さや水分量などの情報を地形コストに統合することで、より精密な経路計画が可能となる。これにより、精密農業の実現に貢献できる。

農業ロボットでは、作業効率が重要である。畝間を効率的に移動し、作業時間を最小化する経路を生成することが求められる。Field D* Hybridの経路短縮効果は、作業効率の向上に直接貢献する。

\subsection{災害対応ロボットへの応用}

災害対応ロボットへの応用も有望である。災害現場では、瓦礫、急斜面、不安定な地盤など、危険な地形が多く存在する。TA*により、これらの危険な領域を回避し、安全な経路を選択することで、ロボットの損傷を防ぎ、救助活動を継続できる。また、経路上の地形情報を事前に把握することで、操縦者の負担を軽減し、より迅速な救助活動が可能となる。これにより、人命救助への貢献が期待できる。

災害対応では、時間が重要である。Field D* Hybridの高速な計算時間（0.175秒）は、迅速な経路計画を可能にし、救助活動の効率を向上させる。

\subsection{惑星探査ローバーへの応用}

さらに、惑星探査ローバーへの応用も考えられる。火星や月などの惑星表面は、クレーター、岩石、砂地など、複雑な地形で構成される。ローバーは、通信遅延が大きいため、自律的に安全な経路を選択する必要がある。TA*を用いることで、危険な地形を回避し、ローバーの損傷リスクを低減できる。また、限られたバッテリー資源を効率的に使用し、探査範囲を拡大できる。これにより、宇宙探査の成功率向上に貢献できる。

惑星探査では、事前の地形情報が限られているため、オンラインでの地形評価と経路計画が重要である。RTAB-Mapと提案手法を統合することで、未知の環境での自律探査が可能となる。

\section{まとめ}

本章では、提案手法の今後の発展方向について述べた。アルゴリズムの拡張として、動的環境対応、マルチモーダル地形評価、機械学習による地形予測などが挙げられる。実環境実験により、提案手法の実用性を検証し、多様な環境での性能を評価することが重要である。また、農業ロボット、災害対応ロボット、惑星探査ローバーなど、様々な分野への応用が期待できる。これらの発展により、提案手法が実用化され、多様なフィールドロボットの性能向上に貢献することが期待される。
